% Part: first-order-logic
% Chapter: completeness
% Section: completeness-theorem

\documentclass[../../../include/open-logic-section]{subfiles}

\begin{document}

\iftag{FOL}
      {\olfileid{fol}{com}{cth}}
      {\olfileid{pl}{com}{cth}}

\olsection{完全性定理}

\begin{explain}
综合上述结果，我们就得到了完全性定理。
\end{explain}

\begin{thm}[完全性定理]
\ollabel{thm:completeness}
设 $\Gamma$ 为一个语句集。若 $\Gamma$ 是一致的，则它是可满足的。
\end{thm}

\begin{proof}
假设 $\Gamma$ 是一致的。\iftag{FOL}{ 根据
  \olref[hen]{lem:henkin}，存在一个饱和的一致集 $\Gamma' \supseteq \Gamma$。}{} 根据 \olref[lin]{lem:lindenbaum}，存在一个一致且完全的 $\Gamma^* \supseteq \iftag{FOL}{\Gamma'}{\Gamma}$。
\iftag{FOL}{
  由于 $\Gamma' \subseteq \Gamma^*$，对于每个
  公式~$!A(x)$，$\Gamma^*$ 都包含一个形如
  \iftag{prvEx}
        {$\lexists[x][!A(x)] \lif !A(c)$}
        {$\lnot\lforall[x][!A(x)] \lif \lnot !A(c)$}
  的语句，因此 $\Gamma^*$ 是饱和的。如果 $\Gamma$
  不包含~$\eq$，则由}{由}
\olref[mod]{lem:truth} 可知，
$\iftag{FOL}{\Sat{M(\Gamma^*)}}{\pSat{v(\Gamma^*)}}{!A}$ 当且仅当 $!A \in
\Gamma^*$。由此特别地，对于所有 $!A \in
\Gamma$，$\iftag{FOL}{\Sat{M(\Gamma^*)}}{\pSat{v(\Gamma^*)}}{!A}$，所以
$\Gamma$ 是可满足的。\iftag{FOL}{
  如果 $\Gamma$ 包含~$\eq$，
  则根据 \olref[ide]{lem:truth}，对于所有语句~$!A$，
  $\Sat{\equivclass{M}{\approx}}{!A}$ 当且仅当 $!A \in \Gamma^*$。特别地，
  对于所有 $!A \in
  \Gamma$，$\Sat{\equivclass{M}{\approx}}{!A}$，所以 $\Gamma$ 是可满足的。}{}
\end{proof}

\begin{cor}[完全性定理，第二形式]
\ollabel{cor:completeness}
对于所有 $\Gamma$ 和语句~$!A$：若 $\Gamma \Entails !A$，则
$\Gamma \Proves !A$。
\end{cor}

\begin{proof}
请注意 \olref{cor:completeness} 和 \olref{thm:completeness} 中的 $\Gamma$ 都是全称量化的。为避免混淆，我们用另一个变元来重述 \olref{thm:completeness}：对于任何语句集~$\Delta$，若 $\Delta$ 是一致的，则它是可满足的。其逆否命题为：若 $\Delta$ 不可满足，则 $\Delta$ 是不一致的。我们将据此证明该推论。

假设 $\Gamma \Entails !A$。那么根据 \olref[syn][sem]{prop:entails-unsat}，$\Gamma \cup \{\lnot !A\}$ 是不可满足的。以 $\Gamma
\cup \{\lnot !A\}$ 作为 $\Delta$，由 \olref{thm:completeness} 的前一形式可得 $\Gamma \cup \{\lnot !A\}$ 是不一致的。根据
\iftag{FOL}{%
  \tagrefs{prfAX/{fol:axd:prv:prop:prov-incons},
    prfSC/{fol:seq:prv:prop:prov-incons},
    prfND/{fol:ntd:prv:prop:prov-incons},
    prfTab/{fol:tab:prv:prop:prov-incons}}}{%
  \tagrefs{prfAX/{pl:axd:prv:prop:prov-incons},
    prfSC/{pl:seq:prv:prop:prov-incons},
    prfND/{pl:ntd:prv:prop:prov-incons},
    prfTab/{pl:tab:prv:prop:prov-incons}}}，
有 $\Gamma \Proves !A$。
\end{proof}

\tagprob{FOL}
\begin{prob}
利用 \olref[fol][com][cth]{cor:completeness} 来证明
\olref[fol][com][cth]{thm:completeness}，从而表明完全性定理的两种表述是等价的。
\end{prob}
\tagendprob

\tagprob{notFOL}
\begin{prob}
利用 \olref[pl][com][cth]{cor:completeness} 来证明
\olref[pl][com][cth]{thm:completeness}，从而表明完全性定理的两种表述是等价的。
\end{prob}
\tagendprob

\tagprob{FOL}
\begin{prob}
要使推导系统具有完全性，其规则必须足够强，以证明每个不可满足的集合都是不一致的。证明完全性时，哪些推导规则是必要的？是否有某些规则在整个证明中完全未被用到？为了回答这些问题，请列一个表或画一个图，说明在导向 \olref[fol][com][cth]{thm:completeness} 证明的各个结果中，分别使用了哪些推导规则。务必留意这些证明中对规则的隐含使用。
\end{prob}
\tagendprob

\tagprob{notFOL}
\begin{prob}
要使推导系统具有完全性，其规则必须足够强，以证明每个不可满足的集合都是不一致的。证明完全性时，哪些推导规则是必要的？是否有某些规则在整个证明中完全未被用到？为了回答这些问题，请列一个表或画一个图，说明在导向 \olref[pl][com][cth]{thm:completeness} 证明的各个结果中，分别使用了哪些推导规则。务必留意这些证明中对规则的隐含使用。
\end{prob}
\tagendprob

\end{document}